% F085｜真实数学图；代表函数不构成抽象解的唯一性证明。
\begin{center}
\begin{tikzpicture}
\begin{axis}[width=12.4cm,height=5.1cm,axis lines=middle,
 xmin=-0.8,xmax=4.5,ymin=-4,ymax=100,
 xtick={0,1,2,3,4},xticklabels=\empty,ytick={0,9,27,81},xlabel={$t$},ylabel={$f(t)$},
 title={例 1 的代表曲线 $3^t$（不声称唯一）},title style={font=\small,yshift=3pt},
 tick label style={font=\small},label style={font=\small},
 ylabel style={at={(axis description cs:0,1)},anchor=south east,xshift=-18pt},
 clip=false,grid=major,grid style={black!8},enlargelimits=false,
 every axis plot/.append style={thick},samples=161]
\addplot[CExplanation,very thick,domain=-.5:4.1] {3^x};
\draw[dashed,black!35] (axis cs:1,0) |- (axis cs:0,3);
\draw[dashed,black!35] (axis cs:2,0) |- (axis cs:0,9);
\draw[dashed,black!35] (axis cs:3,0) |- (axis cs:0,27);
\draw[dashed,black!35] (axis cs:4,0) |- (axis cs:0,81);
\addplot[only marks,mark=*,mark size=2pt,orange!80!black] coordinates {(1,3) (2,9) (3,27) (4,81)};

\node[font=\small,anchor=north,yshift=-3pt] at (axis cs:1,-4) {$1$};
\node[font=\small,anchor=north,yshift=-3pt] at (axis cs:2,-4) {$2$};
\node[font=\small,anchor=north,yshift=-3pt] at (axis cs:3,-4) {$3$};
\node[font=\small,anchor=north,yshift=-3pt] at (axis cs:4,-4) {$4$};

\end{axis}
\end{tikzpicture}
\par\small 输入每加 $1$，高度就乘 $f(1)=3$；四次相乘得到 $f(4)=81$。
\end{center}
